\documentclass{leveldoc}
\DocUses{foundations,deflection_plates,glass_screen,first_coil,first_magnet}
\DocFingerprint{40ca1b3f000d5ca0}

\begin{document}

\section{What the level asks}\label{sec:asks}

\leveldescription

One positive ion is launched along $+x$ from the left edge. Above its line stands a strong
magnet, its moment out of the board, and between the two lies a bar of ferrite: a magnetically
soft insulator of relative permeability $\mu_\mathrm{r} = \val{muR}$. It carries no current,
but in the magnet's field it magnetizes, and a bound magnetic charge appears on its surfaces.
The player places up to three magnets, of either orientation and the listed strengths, inside
the marked region below the launch line. The ion must enter the detector, the rectangle
$\val{xDetLo} \le x \le \val{xDetHi}$, $\val{yDetLo} \le y \le \val{yDetHi}$ (no acceptance
window, no kinetic window, no gate). A flight ends when the ion touches a magnet or the ferrite
(its box plus a contact shell of $\val{contact}$), leaves the world, or enters the detector
(\xsecref{foundations}{sec:board}).

With nothing placed the magnet's field, weakened by the ferrite, still bends the ion upwards,
and it leaves through the right edge above the detector. The reference places a weaker magnet
of the opposite orientation below the path, whose field bends the ion back down into the
detector's left face. Without the ferrite the same magnet would not be enough. This document
derives the ferrite's bound charge, the magnetic twin of the Glass screen's
(\levelref{glass_screen}), solves it on rectangles, and integrates the flights.

\section{Setup}\label{sec:setup}

\setuptable

The ion has charge $q = \val{q}$, mass $m = \val{m}$ and radius $0$; $c = \val{c}$; it starts at
$(\val{xZero}, \val{yZero})$ along $+x$ with $T_0 = \val{Tzero}$, so $\gamma_0 = \val{gammaZero}$
and $p = \val{pZero}$. The level's magnet has the strength $\mu_1 = \val{muOne}$ (field units,
\physref{2.2}) at $(\val{xMagOne}, \val{yMagOne})$; the reference's, $\mu_2 = \val{muTwo}$ at
$(\val{xMagTwo}, \val{yMagTwo})$; each sphere has radius $\val{Rm}$. The ferrite occupies
$\val{xF0} \le x \le \val{xF1}$, $\val{yF0} \le y \le \val{yF1}$, length $\val{fL}$, thickness
$\val{fT}$ and height $\val{fH}$ along $z$, symmetric about the plane of the flight;
$\lambda = (\mu_\mathrm{r} - 1)/(\mu_\mathrm{r} + 1) = \val{lam}$. The world is
$-1 \le x \le \val{xWorld}$, $-1 \le y \le \val{yWorld}$. The reported step is $h = \val{H}$ and
the reported mesh has rectangles of edge $\val{panelFine}$, $\val{nPanels}$ of them on the
ferrite's upper half.

\begin{figure}[ht]
\centering
\begin{tikzpicture}
\begin{axis}[axis equal image, width=0.92\linewidth,
  xmin=-1.5, xmax=32.5, ymin=-2, ymax=22.5,
  xlabel={$x$}, ylabel={$y$}, tick label style={font=\small}]
  \draw[black!40, dashed] (-1,-1) rectangle (\rawval{xWorld},\rawval{yWorld});
  \draw[black!25] (0,0) rectangle (30,20);
  \fill[green!50!black, fill opacity=0.18]
    (\rawval{xDetLo},\rawval{yDetLo}) rectangle (\rawval{xDetHi},\rawval{yDetHi});
  \draw[green!40!black]
    (\rawval{xDetLo},\rawval{yDetLo}) rectangle (\rawval{xDetHi},\rawval{yDetHi});
  \fill[brown!40!white]
    (\rawval{xF0},\rawval{yF0}) rectangle (\rawval{xF1},\rawval{yF1});
  \draw[brown!50!black]
    (\rawval{xF0},\rawval{yF0}) rectangle (\rawval{xF1},\rawval{yF1});
  \fill[black!75] (\rawval{xMagOne},\rawval{yMagOne}) circle[radius=\rawval{Rm}];
  \fill[black!40] (\rawval{xMagTwo},\rawval{yMagTwo}) circle[radius=\rawval{Rm}];
  \addplot[orange!80!black] table {pathbare.dat};
  \addplot[very thick, blue!60!black] table {pathref.dat};
  \fill (\rawval{xZero},\rawval{yZero}) circle[radius=0.08];
\end{axis}
\end{tikzpicture}
\caption{The ferrite (brown), the level's magnet above it (black) and the reference's below the
path (gray); the green box is the detector. Blue: the reference flight, bent down into the
detector's left face. Orange: the empty board, bent up and out through the right edge above
the detector.}
\label{fig:board}
\end{figure}

\section{Principles and model}\label{sec:model}

\begin{principles}
  \item The Lorentz force in the relativistic equation of motion
    \xeqref{foundations}{eq:eom}; the magnetic force does no work
    (\xsecref{first_coil}{sec:nowork}).
  \item The magnets' fields outside their spheres are point dipoles'
    (\xsecref{first_magnet}{sec:dipole}); in three dimensions,
    $\vect{H} = \mu\,(3(\uvec{z}\cdot\uvec{r})\uvec{r} - \uvec{z})/r^3$ (\Jackson{§5.6}).
  \item Magnetostatics without free currents: $\nabla\times\vect{H} = 0$, so
    $\vect{H} = -\nabla\Phi_\mathrm{M}$; $\nabla\cdot\vect{B} = 0$, with $\vect{B} = \vect{H}$
    outside and $\vect{B} = \mu_\mathrm{r}\vect{H}$ in a linear soft ferrite; the normal $\vect{B}$
    is continuous at a surface (\Jackson{§5.8, §5.9}).
  \item The rectangle integral \xeqref{deflection_plates}{eq:rectangle} and the bound-charge
    method of \xsecref{glass_screen}{sec:panels}.
  \item The board's rules, the classical Runge--Kutta method and the cubic Hermite location of an
    event (\xsecref{foundations}{sec:board}, \xsecref{foundations}{sec:numerics}), as in
    \xsecref{deflection_plates}{sec:flight}.
\end{principles}

\begin{assumption}{A linear, static ferrite}\label{as:linear}
$\mu_\mathrm{r}$ is constant: no saturation, no hysteresis, magnetized at once by the magnets
only (the ion's own field is left out). The game makes the same assumptions (\physref{2.7},
\emph{Ferrites}).
\end{assumption}

\begin{assumption}{No radiation}\label{as:radiation}
The radiated energy \xeqref{foundations}{eq:larmor} is not taken from the ion; along the
reference flight it is a fraction $\val{EradFrac}$ of $T_0$.
\end{assumption}

\begin{assumption}{Equal rectangles}\label{as:panels}
The bound magnetic charge is constant on equal rectangles of edge $\val{panelFine}$, its
condition held on average at four Gauss points; the engine holds it on graded triangles
(\physref{2.7}). Doubling the edge moves an end state by $\val{errMesh}$, this mesh's
discretization estimate; halving the step, by $\val{errStep}$.
\end{assumption}

\section{Derivation}\label{sec:derivation}

\subsection{The bound magnetic charge}\label{sec:charge}

Outside the ferrite and inside it there is no current, so $\vect{H} = -\nabla\Phi_\mathrm{M}$
everywhere, and a uniform linear medium has no volume magnetic charge: the ferrite acts through
a surface density $\sigma$, whose field is computed as a charge's. The argument of
\xsecref{glass_screen}{sec:condition}, with $\vect{H}$ for $\vect{E}$ and the continuity of the
normal $\vect{B}$, $\mu_\mathrm{r}H_n^\mathrm{in} = H_n^\mathrm{out}$, for that of the normal
$\vect{D}$, gives
\begin{equation}\label{eq:condition}
  \sigma = \frac{\lambda}{2\pi}\,\bar H_n, \qquad \lambda = \frac{\mu_\mathrm{r} - 1}{\mu_\mathrm{r} + 1},
\end{equation}
with $\bar H_n$ the outward normal field of every source but the surface element itself. With
$\mu_\mathrm{r} = \val{muR}$, $\lambda$ is nearly $1$: the ferrite nearly expels $\vect{H}$ from
itself, drawing the field lines into it.

\subsection{Odd under the mirror}\label{sec:odd}

A moment along $z$ at $z = 0$ has $H_z$ even in $z$ and $H_x$, $H_y$ odd. On the top face,
whose outward normal is $+\uvec{z}$, $\bar H_n = H_z$; on its mirror, the bottom face,
$\bar H_n = -H_z$; on the side faces $\bar H_n$ is an in-plane component, odd in $z$. So
\begin{equation}\label{eq:odd}
  \sigma(x, y, -z) = -\sigma(x, y, z) :
\end{equation}
each rectangle's mirror carries the opposite density. Two consequences: the ferrite's total
magnetic charge is zero by symmetry (no constraint or extra unknown is needed, unlike the glass's
neutrality), and in the plane $z = 0$ the in-plane fields of a rectangle and its mirror cancel,
while their $z$ components add: the ferrite's field in the plane is perpendicular to it, and the
ion stays in the plane. On the reported mesh the upper half carries $\val{qHalf}$, the lower half
the opposite.

\subsection{The panels}\label{sec:panels}

Rectangle $j$ at unit density has the field $\vect{e}_j$ (the gradient of
\xeqref{deflection_plates}{eq:rectangle}); its mirror at unit density $\vect{e}_j^\mathrm{m}$.
Averaging \eqref{eq:condition} over rectangle $i$ at its four Gauss points $\vect{g}_{i,k}$ and
projecting on its outward normal $\uvec{n}_i$,
\begin{equation}\label{eq:system}
  \sigma_i - \frac{\lambda}{2\pi}\sum_j \sigma_j K_{ij}
    = \frac{\lambda}{2\pi}\,\frac14\sum_k \vect{H}_\mathrm{ext}(\vect{g}_{i,k})\cdot\uvec{n}_i,
  \qquad
  K_{ij} = \frac14\sum_k\big(\vect{e}_j - \vect{e}_j^\mathrm{m}\big)(\vect{g}_{i,k})\cdot\uvec{n}_i ,
\end{equation}
with $\vect{H}_\mathrm{ext}$ the magnets' dipole fields, evaluated off the plane at the Gauss
points, and the rectangle's own $\vect{e}_i$ left out of $K_{ii}$ (its normal part averages to
zero over its two sides; its mirror stays). The dense system is solved by LU factorization;
its residual is $\val{solveRes}$. In the plane,
\begin{equation}\label{eq:bz}
  B_z(x, y) = -\sum_{\text{magnets}}\frac{\mu_k}{\abs{\vect{x} - \vect{r}_k}^3}
    + \sum_j \sigma_j\big(e_{j,z} - e^\mathrm{m}_{j,z}\big)(x, y, 0),
\end{equation}
and the ion obeys \xeqref{first_magnet}{eq:eom-dipole} with this $B_z$. The field with nothing
placed, the level's magnet alone against magnet and ferrite:
\begin{center}\small\val{fieldtab}\end{center}
On the launch line under the magnet the ferrite leaves $\val{shieldRatio}$ of the field: it
shields, unlike the glass of \levelref{glass_screen}, because $\lambda \approx 1$ and the magnet's
field crosses the bar's large top and bottom faces, where the bound charge sits and cancels it.

\subsection{The flight}\label{sec:flight}

The state is advanced by \xeqref{foundations}{eq:rk4} with $h = \val{H}$, the ending located on
the cubic Hermite interpolant (thirty-two samples, sixty halvings) with obstacles first, as in
\xsecref{deflection_plates}{sec:flight}. The magnetic force does no work, so $\abs{\vect{p}}$ is
constant: the check. The step's error is $\abs{y_h - y_{2h}}/15$; the mesh estimate is the
change when the rectangles' edge doubles to $\val{panelCoarse}$. Both flights keep their
outcome at $h$, $2h$ and on both meshes.

\section{The reference solution by hand}\label{sec:hand}

\begin{procedure}
  \item\label{st:launch} \emph{The launch.} $\gamma_0 = \val{gammaZero}$, $\vect{p}_0 = (\val{pZero}, 0)$.
  \item\label{st:panels} \emph{The panels.} Build the upper-half rectangles of edge
    $\val{panelFine}$ ($\val{nPanels}$ of them) and solve \eqref{eq:system} with both magnets'
    fields on the right side: residual $\val{solveRes}$.
  \item\label{st:march} \emph{The march.} Advance \xeqref{first_magnet}{eq:eom-dipole} with $B_z$
    from \eqref{eq:bz} by \xeqref{foundations}{eq:rk4}. The state every five time units:
    \begin{center}\small\val{states}\end{center}
  \item\label{st:end} \emph{The arrival.} After $\val{RefSteps}$ steps the ion enters the
    detector through its \val{RefFace} face at $t = \val{RefT}$, at $(\val{RefX}, \val{RefY})$,
    with momentum $(\val{RefPx}, \val{RefPy})$, $\val{RefAng}$ degrees from $+x$.
  \item\label{st:check} \emph{The checks.} $\abs{\vect{p}}$ drifts by $\val{RefPdev}$; the step
    estimate is $\val{errStep}$, the mesh estimate $\val{errMesh}$; the ferrite's sampled
    clearance is $\val{RefDmin}$, outside its shell $\val{contact}$, with at most
    $\val{stepReach}$ travelled in a step.
\end{procedure}

\begin{result}[The reference flight]
The ion enters the detector through its \val{RefFace} face at $(\val{RefX}, \val{RefY})$ at
$t = \val{RefT}$. The level is solved.
\end{result}

\section{Comparison with the game}\label{sec:comparison}

\comparisontable

The tolerance on a state, $\val{tolFlight}$, is ten times the larger of our estimates (the
mesh's, $\val{errMesh}$), fixed before comparing; the game's own preview--verification gap,
$\val{gapGame}$, is below it. The differences, $\val{cmpRefYcDiff}$ in the reference's ordinate
and $\val{cmpBareYcDiff}$ in the empty board's, are below both discretization estimates, ours
(rectangles of $\val{panelFine}$) and the game's (graded triangles of 0.35 of a cell): the two
meshes' errors, not a difference of model. The ferrite converges more slowly than
the glass (Level 22's differences were a thousand times below its estimate): as $\lambda \to 1$
the bound charge's discretization error grows (\physref{2.7}, the cube at $\mu \to \infty$). The
outcomes agree, and the faces the flights meet agree exactly.

\section{Other solutions}\label{sec:others}

The empty board, the second placement in \code{placements.json}: the ion leaves through the
\val{BareFace} edge at $(\val{BareX}, \val{BareY})$ at $t = \val{BareT}$, bent up by the magnet's
residual field. Without the ferrite it would be bent further, to $y = \val{bareAloneY}$ at the
same edge (our run, not compared). The reference's magnet below the path, of the opposite
orientation, adds a field of the other sign there and turns the ion down into the detector;
without the ferrite the same placement ends as \val{noFerriteOut} through the \val{noFerriteFace}
edge at $y = \val{noFerriteY}$, just above the detector: the solution counts on the shield.
A map of solutions would take the game's verdicts, which this document does not reproduce.

\section{What the game computes}\label{sec:engine}

The engine solves the same condition (\physref{2.7}, \emph{Ferrites}): the dielectric
condition with $\vect{H}$ for $\vect{E}$ and $\mu$ for $\varepsilon$, on graded triangles,
averaged per triangle by a six-point rule, each triangle's mirror carrying $-\sigma$, magnetized
by the level's and the player's static magnets and coils. It integrates
\xeqref{first_magnet}{eq:eom-dipole} adaptively (\physref{5}, \physref{6}, \physref{7}).
This document holds the same condition on equal rectangles; the agreement asked of the two is
$\val{tolFlight}$, and they agree within each one's own discretization estimate.

\section{Further reading}\label{sec:reading}

Magnetostatics with permeable media, the magnetic scalar potential and magnetic shielding:
\Jackson{§5.9, §5.12}. The dielectric twin of this calculation: \levelref{glass_screen}; the
magnets' field and the motion in it: \levelref{first_magnet}.

\levelfileappendix

\end{document}
